% Source: https://artofproblemsolving.com/community/c3h2231355_kor_grade_10_sat_msm_level
% Stable identity: aops-topic-2231355-q30
\textbf{$x=1$.} $f(f(1))+f^{-1}(1)=2$ and both terms are at least $1$, so $f^{-1}(1)=1$, i.e. $f(1)=1$.

\textbf{$x=2$.} Since $f(1)=1$, neither $f^{-1}(2)$ nor $f(2)$ equals $1$, and $f(f(2))=1$ would force $f(2)=1$. So both terms of $f(f(2))+f^{-1}(2)=4$ are at least $2$, hence both equal $2$, and $f(2)=2$.

So $f$ maps $\{3,4,5,6,7\}$ onto itself.

\textbf{$x=6$.} Let $a=f(6)$ and $b=f^{-1}(6)$, so $f(b)=6$. Both $a,b\in\{3,4,5,7\}$ (by III, $a\ne6$, and $b=6$ would mean $f(6)=6$). Condition I gives $f(a)+b=12$ with $f(a)\le7$ and $b\le7$, so $b\in\{5,7\}$.

\emph{Case $b=5$.} Then $f(5)=6$, so $f(3)=4$ by II, and $f(a)=7$. Now $a\ne3$ ($f(3)=4$), $a\ne5$ ($f(5)=6$) and $a\ne7$ (then $f(6)=f(7)=7$), so $a=4$, i.e. $f(6)=4=f(3)$: impossible.

\emph{Case $b=7$.} Then $f(7)=6$ and $f(a)=5$. The values of $f$ on $\{3,4,5,6\}$ are $\{3,4,5,7\}$, and II forces $\{f(3),f(5)\}=\{3,7\}$, so $\{f(4),f(6)\}=\{4,5\}$. If $f(6)=5$ then $f(a)=f(5)=5$, contradicting $f(5)\in\{3,7\}$. Hence $f(6)=4$ and $f(4)=5$ (and indeed $f(a)=f(4)=5$).

Both remaining choices $f(3)=3,f(5)=7$ and $f(3)=7,f(5)=3$ satisfy all the conditions (for $x=6$: $f(f(6))+f^{-1}(6)=f(4)+7=12$). In either case
\[f(4)\cdot\bigl(f(6)+f(7)\bigr)=5\cdot(4+6)=\boxed{50}.\]
